\section{Выводы}

Выполнив четвертую лабораторную работу , я глубоко изучил алгоритм поиска подстроки Кнута-Морриса-Пратта и его связь с Z-algorithm, а также разобрался в сдвигах и оптимизациях, которые позволяют достичь линейной сложности

В процессе выполнения лабораторной работы я:

\begin{itemize}
    \item Детально разобрал теорию и практику алгоритма KMP, включая построение функции отказов через промежуточный Z-array и конвертацию в strong prefix function для оптимизированного поиска.

    \item Реализовал эффективный парсер текста, работающий со строками переменной длины и сохраняющий информацию о позициях (строка, позиция слова).

    \item Применил технику потоковой обработки с циклическим буфером для минимизации использования памяти при работе с большими объемами данных.

    \item Научился анализировать результаты бенчмарков и объяснять причины несовпадения теоретических ожиданий с практическими результатами.
\end{itemize}

Реализация демонстрирует правильное понимание алгоритма KMP, применение Z-algorithm для построения таблиц отказов и эффективную работу с потоками данных.

\pagebreak
